%
% Teil 5: Numerische Inversion -- Oszillationsproblem des
% Bromwich-Integrals, Pfaddeformation nach Cauchy, Talbot-Kontur,
% Trapezregel und Konvergenzvergleich
%
\section{Numerische Inversion: Talbots Weg ins Tal des Integranden}
\label{laplace:subsection:talbot}

In der Praxis endet die analytische Inversion schnell in einer Sackgasse.
Viele Bildfunktionen besitzen keine geschlossen darstellbare
Originalfunktion, oder die gefundenen Reihen und Integrale konvergieren
für die interessierenden Zeiten quälend langsam.
Gefragt ist dann eine numerische Auswertung des Bromwich-Integrals, und
der naheliegendste Versuch wäre eine \emph{Quadratur}, also eine
numerische Integration über gewichtete Funktionswerte an endlich vielen
Stützstellen, direkt entlang der Bromwich-Geraden.
Dieser Versuch scheitert spektakulär.
Auf der Geraden $s = \gamma + i\omega$ gilt
$e^{st} = e^{\gamma t} \left( \cos \omega t + i \sin \omega t \right)$;
der Betrag des Exponentialfaktors ist konstant, und der Integrand
oszilliert unablässig, während er nur so langsam abklingt, wie es die
Bildfunktion $F(s)$ selbst vorgibt.
Die Quadratur muss dann riesige, sich fast vollständig auslöschende
Beiträge aufsummieren, und die Genauigkeit wächst mit der Zahl der
Stützstellen nur schleppend.
Abbildung~\ref{laplace:fig:integrand_vergleich} zeigt dieses Verhalten für
das einfachste denkbare Beispiel $F(s) = \frac{1}{s}$.

\begin{figure}
    \centering
    % ============================================================================
% PLATZHALTER -- Bild 5: Entrollter Integrand, Bromwich vs. Talbot
% ============================================================================
% REGIEANWEISUNG FUER DIE GRAFIK-TOOLCHAIN (bevorzugt matplotlib/pgfplots,
% als PDF eingebunden oder direkt pgfplots; Schrift Times 10pt wie Text):
%
% DATENGRUNDLAGE (echte Berechnung, F(s)=1/s, t=1):
% - BROMWICH: s = gamma + i*omega mit gamma=1; Integrand
%   Re[ e^{st} F(s) ] = Re[ e^{t(gamma+i*omega)} / (gamma+i*omega) ]
%   als Funktion des Kurvenparameters omega in [-40, 40].
%   -> oszillierende Kurve mit nur ~1/|omega| abfallender Amplitude.
% - TALBOT: Kontur s(theta) = sigma + lambda(theta*cot(theta) + i*beta*theta)
%   mit z.B. sigma=0, lambda=2, beta=1, theta in (-pi, pi);
%   (beta ist der Formparameter, im Paper-Text s_beta; in der
%   Originalliteratur oft nu genannt -- hier NICHT nu beschriften!)
%   Integrand Re[ e^{t s(theta)} F(s(theta)) s'(theta) / (2*pi*i) ]
%   als Funktion von theta.
%   -> glatte, glockenfoermige Kurve, an den Raendern theta -> +-pi
%      praktisch null (superexponentieller Abfall).
%
% DARSTELLUNG:
% - ZWEI Panels nebeneinander (gleiche Hoehe), oder ein Panel mit zwei
%   Kurven, falls die Skalen vertraeglich sind; bevorzugt ZWEI PANELS:
%   * links:  "Bromwich-Gerade", Kurve in ROT (plotred), x-Achse
%     "$\omega$", y-Achse "Integrand".
%   * rechts: "Talbot-Kontur", Kurve in BLAU (plotblue), x-Achse
%     "$\theta$", y-Achse identisch skaliert beschriftet.
% - Beide y-Achsen mit derselben Einheit; falls die Amplituden stark
%   differieren, y-Limits pro Panel waehlen, aber im Caption-Text bleibt
%   die Aussage qualitativ.
% - Dezentes Gitter (grau, duenn), keine Titelzeile im Bild (Caption
%   uebernimmt das), Legenden nur falls beide Kurven in einem Panel.
% - Auf der Talbot-Kurve zusaetzlich N=12 Stuetzstellen der Trapezregel
%   als Punkte markieren (blaue Punkte), um zu zeigen, wie wenige
%   Auswertungen genuegen.
%
% GROESSE: Gesamtbreite ca. 12cm, Hoehe ca. 4.5cm.
% ============================================================================
\scalebox{0.8}{
\begin{tikzpicture}

% ==============================
% 1. FARBEN DEFINIEREN
% ==============================
\definecolor{plotblue}{RGB}{0, 0, 200}
\definecolor{plotred}{RGB}{220, 30, 30}
\definecolor{plotgray}{RGB}{140, 140, 140}
\definecolor{plotbg}{RGB}{244, 244, 252}

% ==============================
% 2. PGFPLOTS SETUP
% ==============================
\pgfplotsset{table/search path={.,Latex/images,images,papers/laplace/images}}

% ==============================
% 3. BROMWICH PANEL (Links)
% ==============================
\begin{axis}[
    name=ax1,
    width=7cm,
    height=5.5cm,
    grid=major,
    grid style={plotgray!30, thin},
    axis lines=left,
    axis line style={thick},
    tick label style={font=\small},
    label style={font=\normalsize},
    enlarge x limits=false,
    xlabel={$\omega$},
    ylabel={Integrand},
    xmin=-40, xmax=40,
    ymin=-0.8, ymax=3.0,
    title={\textbf{Bromwich-Gerade}},
    title style={font=\normalsize, yshift=-1.5ex}
]
\addplot[
    thick, 
    plotred,
    ] table[x=omega, y=integrand, col sep=comma] {papers/laplace/images/data_integrand_bromwich.csv};
\end{axis}

% ==============================
% 3. TALBOT PANEL (Rechts)
% ==============================
\begin{axis}[
    name=ax2,
    at={(ax1.outer east)},
    anchor=outer west,
    xshift=1.5cm,
    width=7cm,
    height=5.5cm,
    grid=major,
    grid style={plotgray!30, thin},
    axis lines=left,
    axis line style={thick},
    tick label style={font=\small},
    label style={font=\normalsize},
    enlarge x limits=false,
    xlabel={$\theta$},
    xmin=-3.1415, xmax=3.1415,
    ymin=-0.3, ymax=1.3,
    xtick={-3.1415, -1.5708, 0, 1.5708, 3.1415},
    xticklabels={$-\pi$, $-\pi/2$, $0$, $\pi/2$, $\pi$},
    title={\textbf{Talbot-Kontur}},
    title style={font=\normalsize, yshift=-1.5ex}
]
\addplot[
    thick, 
    plotblue,
    ] table[x=theta, y=integrand, col sep=comma] {papers/laplace/images/data_integrand_talbot.csv};

\addplot[
    only marks,
    mark=*,
    mark size=1.5pt,
    plotblue,
    ] table[x=theta, y=integrand, col sep=comma] {papers/laplace/images/data_integrand_talbot_pts.csv};
\end{axis}

\end{tikzpicture}
}

    \caption{Der entlang der Bogenlänge entrollte Realteil des
    Inversionsintegranden für $F(s) = \frac{1}{s}$ bei festem $t$.
    Auf der Bromwich-Geraden (rot) oszilliert die Funktion mit nur
    langsam abfallender Amplitude, auf der Talbot-Kontur (blau) ist sie
    eine glatte, rasch abklingende Kurve, die sich mit wenigen
    Stützstellen präzise integrieren lässt.}
    \label{laplace:fig:integrand_vergleich}
\end{figure}

Die Rettung liegt in einer Freiheit, die wir längst besitzen, aber bisher
nur zum Schliessen der Kontur genutzt haben.
Nach dem Cauchyschen Integralsatz dürfen wir den Integrationsweg beliebig
deformieren, solange wir dabei keine Singularität von $F(s)$ überstreichen
und $F(s)$ im überstrichenen Gebiet gleichmässig gegen null abfällt; der
Wert des Integrals bleibt dabei exakt gleich.
Diese Idee ist alt, denn schon \cite{laplace:carslaw1959} verformt den
Bromwich-Pfad zu Kurven, die links im Unendlichen beginnen und enden, um
die Konvergenz der Integrale nachzuweisen.
Talbot hat daraus in seiner wegweisenden Arbeit
\cite{laplace:talbot1979} ein numerisches Verfahren gemacht.
Sein Gedanke:
Man biege die beiden Enden des Pfades in die linke Halbebene, sodass
$\text{Re}(s) \to -\infty$ läuft.
Dort fällt der Faktor $|e^{st}| = e^{t\,\text{Re}(s)}$ exponentiell ab, der
Integrand wird winzig, und die Oszillationen verlieren jede Bedeutung.
Anschaulich verlässt der Pfad den Bergkamm der Betragsfläche
$|e^{st}F(s)|$ und steigt ins \emph{Tal} der linken Halbebene hinab.
In der logarithmischen Farbskala von
Abbildung~\ref{laplace:fig:betrags_landschaft} wird dieser Abstieg
unmittelbar sichtbar:
Die Bromwich-Gerade verharrt in den hellen Zonen konstanter Höhe,
während die Talbot-Kontur beidseitig in die dunklen Bereiche des
exponentiell kleinen Betrags eintaucht.

\begin{figure}
    \centering
    % ============================================================================
% PLATZHALTER -- Bild 6: Heatmap |e^{st} F(s)| mit Bromwich- und
%                        Talbot-Pfad
% ============================================================================
% REGIEANWEISUNG FUER DIE GRAFIK-TOOLCHAIN (matplotlib empfohlen,
% Export als PDF; Schrift Times 10pt; KEINE transparenten PNG-Teile):
%
% DATENGRUNDLAGE (F(s)=1/s, t=1):
% - Gitter der s-Ebene: Re(s) in [-15, 5], Im(s) in [-12, 12],
%   Aufloesung >= 800x800.
% - Farbwert: log10( |e^{s*t} / s| ), geclippt auf z.B. [-8, 2].
%
% DARSTELLUNG:
% - 2D-CONTOUR-HEATMAP (pcolormesh oder contourf mit vielen Stufen),
%   sequentielle Colormap (z.B. viridis); COLORBAR rechts mit
%   Beschriftung "$\log_{10}|e^{st}F(s)|$".
% - Zusaetzlich duenne weisse Konturlinien bei ganzzahligen log10-Stufen,
%   damit die "Landschaft" (Bergkamm rechts, Tal links) plastisch wird.
% - SINGULARITAET s=0: kleines weisses/rotes Kreuz mit Label "Pol".
% - BROMWICH-PFAD: vertikale ROTE Linie bei Re(s)=1 ueber die volle
%   Bildhoehe, Pfeilspitze nach oben, Label "$\Gamma_B$" (rot).
% - TALBOT-PFAD: BLAUE Kurve s(theta)=sigma+lambda(theta*cot(theta)
%   + i*beta*theta) mit denselben Parametern wie in Bild 5 (sigma=0,
%   lambda=2, beta=1; Formparameter beta, nicht nu beschriften),
%   Pfeilspitzen in Durchlaufrichtung (von unten nach
%   oben), Label "$\Gamma_T$" (blau).
%   Die Kurve muss erkennbar ins dunkle "Tal" links abtauchen und den
%   Pol rechts umschliessen.
% - Textannotation "Tal: $|e^{st}|$ exponentiell klein" im linken
%   Bildbereich (weiss, small); optional "Bergkamm" rechts.
% - Achsenbeschriftung "Re(s)", "Im(s)"; kein Bildtitel.
%
% ALTERNATIVE (falls 3D gewuenscht): 3D-Surface derselben Daten mit
% Blick von schraeg oben, Pfade auf die Flaeche projiziert; 2D-Version
% ist fuer den Druck aber besser lesbar und wird bevorzugt.
%
% GROESSE: ca. 11cm x 8cm inkl. Colorbar, 300dpi falls Rasterexport.
% ============================================================================
\scalebox{0.85}{
\begin{tikzpicture}

% ==============================
% 1. FARBEN DEFINIEREN
% ==============================
\definecolor{plotblue}{RGB}{0, 0, 200}
\definecolor{plotred}{RGB}{220, 30, 30}
\definecolor{plotgray}{RGB}{140, 140, 140}
\definecolor{plotbg}{RGB}{244, 244, 252}

% PGFPLOTS SETUP
\pgfplotsset{table/search path={.,Latex/images,images,papers/laplace/images}}

\begin{axis}[
    width=11cm,
    height=9.5cm,
    enlargelimits=false,
    axis on top,
    xlabel={$\operatorname{Re}(s)$},
    ylabel={$\operatorname{Im}(s)$},
    xmin=-15, xmax=5,
    ymin=-12, ymax=12,
    tick label style={font=\small},
    label style={font=\normalsize},
    colorbar,
    colormap/viridis,
    point meta min=-8,
    point meta max=2,
    colorbar style={
        ylabel={$\log_{10}|e^{st}F(s)|$},
        ytick={-8,-6,-4,-2,0,2},
        ylabel style={font=\normalsize},
        tick label style={font=\small}
    }
]

% Hintergrund-Heatmap (Bild, darauf projiziert)
\addplot graphics [xmin=-15, xmax=5, ymin=-12, ymax=12] {papers/laplace/images/heatmap_bg.png};

% Singularitaet s=0 (Pol)
\draw[thick, white, line cap=round] (axis cs:-0.4,-0.4) -- (axis cs:0.4,0.4) (axis cs:-0.4,0.4) -- (axis cs:0.4,-0.4);
\draw[thick, plotred, line cap=round] (axis cs:-0.2,-0.2) -- (axis cs:0.2,0.2) (axis cs:-0.2,0.2) -- (axis cs:0.2,-0.2);
\node[white, left, font=\small, xshift=-2pt] at (axis cs:-0.2, 0) {\textbf{Pol}};

% Bromwich-Gerade (rot)
\draw[thick, plotred, decoration={markings, mark=at position 0.65 with {\arrow{Stealth[length=3mm, width=2mm]}}}, postaction={decorate}] 
    (axis cs:1, -12) -- (axis cs:1, 12);
\node[plotred, right, font=\large] at (axis cs:1.2, 9) {$\Gamma_B$};

% Talbot-Kontur (blau)
\addplot[
    thick, plotblue,
    decoration={markings, mark=at position 0.35 with {\arrow{Stealth[length=3mm, width=2mm]}}},
    decoration={markings, mark=at position 0.65 with {\arrow{Stealth[length=3mm, width=2mm]}}},
    postaction={decorate}
    ] table[col sep=comma, x=re, y=im] {papers/laplace/images/talbot_contour_heatmap.csv};
\node[plotblue, left, font=\large] at (axis cs:-3, 6) {$\Gamma_T$};

% Annotation Tal removed

% ==============================
% 2. LEGENDE
% ==============================
\filldraw[fill=white, draw=black, thick, rounded corners=3pt] (axis cs:-14.5, 11.5) rectangle (axis cs:-2.5, 8.5);
\draw[thick, plotred, decoration={markings, mark=at position 0.55 with {\arrow{Stealth[length=3mm, width=2mm]}}}, postaction={decorate}] (axis cs:-14.0, 10.5) -- (axis cs:-12.5, 10.5);
\node[right] at (axis cs:-12.0, 10.5) {Bromwich-Gerade $\Gamma_B$};
\draw[thick, plotblue, decoration={markings, mark=at position 0.55 with {\arrow{Stealth[length=3mm, width=2mm]}}}, postaction={decorate}] (axis cs:-14.0, 9.5) -- (axis cs:-12.5, 9.5);
\node[right] at (axis cs:-12.0, 9.5) {Talbot-Kontur $\Gamma_T$};

\end{axis}
\end{tikzpicture}
}

    \caption{Betragslandschaft $|e^{st}F(s)|$ über der komplexen
    $s$-Ebene für $F(s) = \frac{1}{s}$ zur Zeit $t = 1$ in
    logarithmischer Farbskala.
    Die Bromwich-Gerade (rot) verläuft auf konstanter Höhe durch das
    oszillierende Gebiet, während die Talbot-Kontur (blau) die
    Singularität umschliesst und beidseitig ins exponentiell tiefe Tal
    der linken Halbebene abtaucht.}
    \label{laplace:fig:betrags_landschaft}
\end{figure}

Konkret verwendet Talbot die Kontur
\begin{equation}
    s(\theta)
    =
    \sigma + \lambda \, s_{\beta}(\theta),
    \qquad
    s_{\beta}(\theta)
    =
    \theta \cot \theta + i \beta \theta,
    \qquad
    -\pi < \theta < \pi,
    \label{laplace:eq:talbot_kontur}
\end{equation}
mit reellen Parametern $\lambda > 0$, $\beta > 0$ und einer Verschiebung
$\sigma$; für den Formparameter schreibt die Originalliteratur meist
$\nu$, was wir hier vermeiden, weil dieser Buchstabe in
Abschnitt~\ref{laplace:subsection:endlicher_stab} bereits als Polindex
vergeben ist.
Die Verschiebung $\sigma$ übernimmt dabei die Rolle der
Bromwich-Konstanten $\gamma$ aus
Gleichung~\eqref{laplace:eq:bromwich_integral}; die Umbenennung folgt der
numerischen Literatur und bezeichnet denselben Rechtsversatz des
Integrationsweges.
Für $\theta \to \pm\pi$ strebt $\theta \cot\theta \to -\infty$, die Kurve
öffnet sich also wie gewünscht nach links und legt sich um die negative
reelle Achse herum; sämtliche Singularitäten unserer
Wärmeleitungsbeispiele, ob Polkette oder Verzweigungsschnitt, werden von
ihr umschlossen.
Setzt man die Parametrisierung in das Bromwich-Integral ein, so entsteht
das Integral über ein endliches Intervall
\begin{equation}
    f(t)
    =
    \frac{\lambda e^{\sigma t}}{2\pi i}
    \int_{-\pi}^{\pi}
    e^{\lambda t s_{\beta}(\theta)}
    \,
    F\bigl(\sigma + \lambda s_{\beta}(\theta)\bigr)
    \,
    s_{\beta}'(\theta)
    \, d\theta.
    \label{laplace:eq:talbot_integral}
\end{equation}
Der Integrand ist analytisch in $\theta$ und verschwindet an beiden
Intervallenden samt allen Ableitungen exponentiell.
Für genau solche Integranden ist die einfachste aller Quadraturformeln,
die Trapezregel, nicht etwa grob, sondern \emph{geometrisch konvergent}:
Jede zusätzliche Stützstelle multipliziert den Fehler mit einem festen
Faktor kleiner als eins \cite{laplace:trefethen2006}.
Der Grund dafür steckt in der Euler-Maclaurin-Formel, nach der der
Trapezfehler aus Randtermen mit den Ableitungen des Integranden an den
Intervallenden besteht:
Da unser Integrand bei $\theta = \pm\pi$ samt allen Ableitungen
abstirbt, verschwinden diese Randterme allesamt, und übrig bleibt nur
ein exponentiell kleiner Rest.
Mit $n$ äquidistanten Knoten $\theta_j = \frac{j\pi}{n}$ und unter
Ausnutzung der Symmetrie der Kontur zur reellen Achse ergibt sich die
Näherung
\begin{equation}
    \tilde{f}(t)
    =
    \frac{\lambda e^{\sigma t}}{n}
    \sum_{j=0}^{n-1}
    \operatorname{Re}
    \left[
        e^{\lambda t s_{\beta}(\theta_j)}
        \,
        F\bigl(\sigma + \lambda s_{\beta}(\theta_j)\bigr)
        \,
        \frac{s_{\beta}'(\theta_j)}{i}
    \right],
    \label{laplace:eq:talbot_summe}
\end{equation}
wobei der Summand für $j = 0$ mit dem Faktor $\frac{1}{2}$ gewichtet wird
\cite{laplace:murli1990}.
Der Randknoten $j = n$, also $\theta = \pi$, fehlt in der Summe nicht aus
Nachlässigkeit:
Dort strebt $\text{Re}\,s_{\beta}(\theta) \to -\infty$, und der Integrand
verschwindet exakt, sodass sein Beitrag zur Trapezsumme entfällt.
Der Realteil-Operator und die halbierten Summationsgrenzen sind dabei
kein Formalismus, sondern gezielte Ausnutzung der Kontursymmetrie:
Weil die Kontur spiegelsymmetrisch zur reellen Achse verläuft und für
reelles $f(t)$ die Beziehung $F(\bar{s}) = \overline{F(s)}$ gilt, liefern
die Stützstellen mit negativem $\theta_j$ exakt die konjugierten Beiträge
ihrer Spiegelpartner.
Statt über alle $2n$ Knoten des vollen Intervalls $(-\pi, \pi)$ zu
summieren, genügen daher die $n$ Knoten mit $\theta_j \geq 0$; die
Paarbildung steckt im Realteil, und der Rechenaufwand halbiert sich.
Die gesamte Inversion schrumpft damit auf eine Handvoll Auswertungen der
Bildfunktion an komplexen Stützstellen zusammen.

Die Parameter sind dabei nicht frei wählbar, sondern müssen zum Problem
passen.
Die Verschiebung $\sigma$ wird rechts der Realteile aller Singularitäten
gewählt, damit die Kontur die Singularitäten tatsächlich umschliesst.
Entscheidend ist die Skalierung von $\lambda$:
Damit die geometrische Konvergenz greift, muss die Kontur mit der
Stützstellenzahl wachsen und mit der Inversionszeit schrumpfen, also
$\lambda \propto \frac{N}{t}$, wobei $N = 2n$ die Gesamtzahl der
Stützstellen auf der vollen Kontur bezeichnet, von denen
Gleichung~\eqref{laplace:eq:talbot_summe} dank der Symmetrie nur die
Hälfte auswertet.
In der von Weideman optimierten Form lautet die Kontur beispielsweise
\begin{equation}
    s(\theta)
    =
    \frac{N}{t}
    \bigl(
        0{,}5017 \, \theta \cot(0{,}6407\,\theta)
        - 0{,}6122
        + 0{,}2645 \, i\theta
    \bigr),
    \label{laplace:eq:weideman_kontur}
\end{equation}
mit der die Rate $O(3{,}89^{-N})$ erreicht wird
\cite{laplace:weideman2006}.
Man beachte, dass hier die \emph{gesamte} Kontur -- Verschiebung und
Skalierung gleichermassen -- mit dem Faktor $\frac{N}{t}$ wächst.
Genau diese Kopplung von Konturgrösse, Stützstellenzahl und
Inversionszeit ist der Kern der Parameterwahl und der Garant für die
geometrische Konvergenz; mit einer festen, von $N$ und $t$ unabhängigen
Kontur ginge sie verloren.
Die Kopplung hat aber auch eine praktische Konsequenz:
Die Stützstellen hängen von $t$ ab, sodass für jeden neuen Zeitpunkt die
Bildfunktion an neuen Punkten ausgewertet werden muss
\cite{laplace:trefethen2006}.

Wie gut das funktioniert, ist verblüffend.
Bereits mit $n = 11$ Stützstellen erreicht Talbots Verfahren typischerweise
Fehler der Ordnung $10^{-6}$, mit $n = 20$ etwa $10^{-11}$ und mit
$n = 35$ in doppelter Genauigkeit rund $10^{-20}$
\cite{laplace:talbot1979}.
Die sorgfältige Software-Implementierung von Murli und Rizzardi bestätigt
diese Zahlen und zeigt in direkten Vergleichen, dass Fourier-basierte
Verfahren für dieselbe Genauigkeit ein Vielfaches an Funktionsauswertungen
benötigen \cite{laplace:murli1990}.
Die moderne Analyse von Trefethen, Weideman und Schmelzer ordnet die
Talbot-Kontur in die Familie der \emph{Hankel-Konturen} ein.
Der Begriff ist dabei ein topologischer Überbegriff:
Er bezeichnet jede Kurve, die aus dem dritten Quadranten von
$\text{Re}(s) = -\infty$ kommend die negative reelle Achse umschlingt
und im zweiten Quadranten wieder nach $\text{Re}(s) = -\infty$
verschwindet -- die Schlüsselloch-Kontur aus
Abschnitt~\ref{laplace:subsection:verzweigung} gehört ebenso dazu wie
Parabeln, Hyperbeln oder Talbots Kotangens-Kontur.
Innerhalb dieser Familie beziffert die Analyse die Konvergenzraten mit
optimierten Parametern:
Parabeln erreichen $O(2{,}85^{-N})$, Hyperbeln $O(3{,}20^{-N})$ und
Talbots Kotangens-Konturen $O(3{,}89^{-N})$ bei $N$ Stützstellen
\cite{laplace:trefethen2006}; die zugrunde liegende Parameteroptimierung
stammt aus \cite{laplace:weideman2006}.
Diese Raten sind keine Eigenheit des Testbeispiels $F(s) = \frac{1}{s}$,
sondern theoretische Schranken, die für die gesamte Klasse von
Bildfunktionen mit Singularitäten auf oder nahe der negativen reellen
Achse gelten -- also für alle unsere Wärmeleitungsprobleme, ob Polkette
oder Verzweigungsschnitt.
Mit $N = 32$ ist damit Maschinengenauigkeit erreicht.
Abbildung~\ref{laplace:fig:konvergenz_vergleich} stellt dieses
geometrische Konvergenzverhalten der langsamen Konvergenz auf der
Bromwich-Geraden gegenüber.

\begin{figure}
    \centering
    % ============================================================================
% Bild 7: Konvergenzvergleich (Fehler vs. Stuetzstellen), nativ in pgfplots
% ============================================================================
% DATEN: data_konvergenz.csv, erzeugt von generate_konvergenz_data.py
% (Testproblem F(s) = 1/s, exakt f(t) = 1, t = 1.5).
%
%   * BROMWICH: Mittelpunktsregel auf s = 1 + i*omega, |omega| <= 2*sqrt(N).
%     Der Abschneidefehler ~ e^{gamma t}/(pi*t*Omega) ~ N^(-1/2) dominiert
%     (der Aliasing-Anteil faellt wie e^{-(pi/2)*sqrt(N)} weg):
%     langsame, leicht oszillierende algebraische Konvergenz.
%   * TALBOT: Mittelpunktsregel auf der Weideman-Kontur (Trefethen/
%     Weideman/Schmelzer 2006, Gl. (3.3), skaliert mit N/t).
%     Geometrische Konvergenz O(3.89^-N) bis ~3e-14 bei N ~ 28; DANACH
%     steigt der Fehler wieder an: der Konturscheitel liegt bei
%     Re(st) ~ 0.171*N, also verstaerkt der Faktor e^{st} den Rundungs-
%     fehler wie eps*e^{0.171*N} (Talbots Fehlerterm E_r) -- bei N = 100
%     sind das ~5e-9. Dieses V-foermige Verhalten ist ECHT und wird
%     bewusst gezeigt (Annotation "Rundungsfehler").
%
% DESIGN-ENTSCHEIDUNGEN:
%   * table/search path deckt beide Build-Kontexte ab: Buch-Build
%     (cwd = Latex/, Pfad images/...) und test_render.py (cwd =
%     Projektwurzel, Pfad Latex/images/...).
%   * Referenzgerade 0.5*3.89^(-N) nur bis N = 26 gezeichnet -- ab N = 28
%     ist sie in den Daten auf 1e-16 gesaettigt und wuerde sonst als
%     horizontale Linie auf dem Boden weiterlaufen.
%   * Umlaute als \"u-Escapes (test_render.py liest die Datei ohne
%     explizites Encoding; rohe UTF-8-Bytes wuerden verstuemmelt).
%   * 17 Dekaden (1e-17..1e0), Ticks alle 4 Dekaden; Maschinengenauigkeit
%     als gepunktete Linie bei eps ~ 1e-16 knapp ueber dem Achsenboden.
% ============================================================================
\begin{tikzpicture}

% ==============================
% 1. FARBEN DEFINIEREN
% ==============================
\definecolor{plotblue}{RGB}{0, 0, 200}
\definecolor{plotred}{RGB}{220, 30, 30}
\definecolor{plotgray}{RGB}{140, 140, 140}
\definecolor{plotbg}{RGB}{244, 244, 252}

% CSV in beiden Build-Kontexten auffindbar machen.
% REIHENFOLGE WICHTIG: Latex/images VOR images, denn im Werkstatt-Root
% existiert ein aelterer images/-Ordner mit veralteten Daten, der sonst
% die frische CSV verschatten wuerde. Beim Buch-Build (cwd = Latex/)
% existiert Latex/images nicht und es greift automatisch images/.
\pgfplotsset{table/search path={.,Latex/images,images,papers/laplace/images}}

\begin{semilogyaxis}[
    width=10.5cm,
    height=7.5cm,
    xlabel={Anzahl der St\"utzstellen $N$},
    ylabel={Absoluter Fehler $|f_{\text{num}} - 1|$},
    xmin=0, xmax=104,
    ymin=1e-17, ymax=1e0,
    xtick={0,20,40,60,80,100},
    ytick={1e0,1e-4,1e-8,1e-12,1e-16},
    grid=both,
    grid style={line width=.1pt, draw=gray!20},
    major grid style={line width=.2pt, draw=gray!50},
    legend cell align={left},
    legend style={at={(0.97,0.82)}, anchor=north east,
        nodes={scale=0.9, transform shape}, draw=black!30,
        rounded corners=2pt, inner xsep=4pt, inner ysep=3pt},
    tick label style={font=\footnotesize},
    label style={font=\small}
]

    % Maschinengenauigkeit (eps ~ 1e-16)
    \draw[black, thick, densely dotted]
        (axis cs:0, 1e-16) -- (axis cs:104, 1e-16);
    \node[black, above right, font=\footnotesize]
        at (axis cs:2, 1.3e-16) {Maschinengenauigkeit};

    % Referenzgerade O(3.89^-N), nur bis N = 24: danach ist sie in den
    % Daten gesaettigt UND wuerde unten in den Schriftzug
    % "Maschinengenauigkeit" hineinlaufen
    \addplot[plotgray, thick, dashed, forget plot]
        table [col sep=comma, x=N, y=err_ref,
               restrict x to domain=4:24] {papers/laplace/images/data_konvergenz.csv};
    % Label rechts oberhalb der Referenzgeraden im leeren Keil zwischen
    % Bromwich-Kurve (oben) und Talbot-Abfall (unten); die fruehere
    % Position am V-Boden (27, 1.5e-15) kollidierte optisch mit den
    % blauen Quadraten bei N = 24-28 und dem Schriftzug
    % "Maschinengenauigkeit"
    \node[plotgray, anchor=west, font=\footnotesize]
        at (axis cs:19, 2e-9) {$\mathcal{O}(3{,}89^{-N})$};

    % Bromwich-Gerade: langsame algebraische Konvergenz (rot, Kreise)
    \addplot[
        thick, plotred,
        mark=o, mark options={solid, scale=0.8}
    ] table [col sep=comma, x=N, y=err_bromwich] {papers/laplace/images/data_konvergenz.csv};
    \addlegendentry{Bromwich-Gerade}

    % Talbot-Kontur: geometrischer Abfall, dann Rundungsfehler-Anstieg
    % (blau, Quadrate)
    \addplot[
        thick, plotblue,
        mark=square, mark options={solid, scale=0.8}
    ] table [col sep=comma, x=N, y=err_talbot] {papers/laplace/images/data_konvergenz.csv};
    \addlegendentry{Talbot-Kontur}

    % Annotation am wieder ansteigenden Ast: eps * e^{0.171 N}
    % (etwas oberhalb der Kurve, damit die blauen Marker bei N ~ 76-84
    % nicht in den Schriftzug laufen)
    \node[plotgray, anchor=south, font=\footnotesize]
        at (axis cs:70, 8e-10) {Rundungsfehler};

\end{semilogyaxis}
\end{tikzpicture}

    \caption{Konvergenzvergleich der numerischen Inversion für
    $F(s) = \frac{1}{s}$ bei festem $t$:
    absoluter Fehler gegen die Anzahl der Stützstellen $N$ in
    halblogarithmischer Darstellung.
    Die Quadratur auf der Bromwich-Geraden (rot) verbessert sich nur
    langsam, während die Trapezregel auf der Talbot-Kontur (blau)
    geometrisch mit der Rate $O(3{,}89^{-N})$ konvergiert und um
    $N \approx 30$ nahe an die Maschinengenauigkeit herankommt.
    Für noch grössere $N$ steigt der Fehler wieder leicht an, weil der
    Faktor $e^{st}$ am Konturscheitel den Rundungsfehler wie
    $\varepsilon \, e^{0{,}17 N}$ verstärkt.}
    \label{laplace:fig:konvergenz_vergleich}
\end{figure}

Zwei Einschränkungen gehören der Vollständigkeit halber dazu.
Erstens muss $|F(s)|$ für $|s| \to \infty$ gleichmässig abfallen, damit
die Pfaddeformation erlaubt ist.
Zweitens müssen die Imaginärteile aller Singularitäten beschränkt sein,
denn die Kontur kann nur Singularitäten umschliessen, die nicht beliebig
weit nach oben oder unten ausweichen \cite{laplace:talbot1979}.
Und selbst innerhalb dieser Grenzen kostet jede Abweichung von der
reellen Achse:
Liegen Singularitäten mit grossem Imaginärteil $q_d$ vor, muss sich die
Kontur weit öffnen, um sie zu umschliessen, und die benötigte
Stützstellenzahl wächst proportional zum Produkt $v = q_d \, t$
\cite{laplace:talbot1979}.
Da die Lage der Singularitäten $q_d$ vom Problem fest vorgegeben ist,
wächst $v$ mit fortschreitender physikalischer Zeit $t$ zwangsläufig an
-- bei Schwingungsproblemen verliert das Verfahren seine Effizienz also
nicht durch einen behebbaren Defekt, sondern durch die Zeit selbst.
Unsere Wärmeleitungsprobleme mit ihren Polen und Schnitten auf der
negativen reellen Achse sind damit der Idealfall des Verfahrens, während
etwa die Schwingungsgleichung mit ihrer unendlichen Polkette auf der
imaginären Achse ausserhalb seines Anwendungsbereichs liegt.
Innerhalb dieses Bereichs ist die Methode jedoch so leistungsfähig, dass
sie heute weit über die Laplace-Inversion hinaus eingesetzt wird, etwa zur
Berechnung von Matrix-Exponentialfunktionen und zur Zeitintegration
parabolischer partieller Differentialgleichungen
\cite{laplace:trefethen2006}.

Damit schliesst sich der Bogen dieser Arbeit.
Aus der Korrespondenztabelle wurde das Bromwich-Integral, aus dem Integral
eine Residuensumme, aus dem Umgang mit Verzweigungsschnitten und
Polketten ein Verständnis dafür, wie die Geometrie des physikalischen
Problems die Singularitäten in der $s$-Ebene formt.
Die Talbot-Kontur schliesslich zeigt, dass dieselbe Funktionentheorie, die
die analytische Inversion trägt, auch das effizienteste numerische
Verfahren begründet:
Man muss das Bromwich-Integral nicht dort auswerten, wo es definiert
wurde, sondern darf es dorthin verschieben, wo es am schnellsten
verschwindet.
